% ============================================================
% RELATED WORK SECTION
% ============================================================

\section{Related Work}\label{sec:related}

Our work builds on 3 research areas: reinforcement learning for sepsis treatment, offline RL algorithms, and interpretability methods for RL policies.

% ============================================================
\subsection{Reinforcement Learning for Sepsis Treatment}\label{sec:related:sepsis}

Raghu et al. \citeyearpar{raghu2017sepsis_drl} pioneered deep RL for sepsis treatment using the MIMIC-III database, formulating treatment as a discrete-action MDP with a $5 \times 5$ action grid for IV fluid and vasopressor dosing. Their work established the gym-sepsis simulation environment we use in this study.

Komorowski et al. \citeyearpar{komorowski2018ai_clinician} developed the AI Clinician, a landmark study in clinical RL that trained a sepsis treatment policy on 96,156 ICU patient records from MIMIC-III and the eICU database. Using fitted Q-iteration with K-nearest neighbors value function approximation, they clustered the continuous 48-dimensional state space into 750 discrete states and optimized a policy over a discrete action space representing IV fluid and vasopressor dosing combinations. The AI Clinician achieved 98.1\% survival in retrospective off-policy evaluation, compared to 86.3\% for observed clinician actions, representing a substantial performance gain. Komorowski et al. demonstrated that the learned policy aligned qualitatively with clinical guidelines by visualizing median action recommendations across patient severity states, showing that the AI Clinician escalated vasopressor dosing for patients with low blood pressure and high lactate, consistent with Surviving Sepsis Campaign recommendations \citep{rhodes2017ssc}.

Despite the AI Clinician's strong performance and qualitative alignment with medical knowledge, Komorowski et al. did not provide quantitative feature-attribution analysis to measure how strongly individual physiological variables such as blood pressure, lactate, and SOFA score drive policy decisions. Their interpretability assessment relied on visualizing aggregate action distributions across discretized state clusters, revealing what the policy recommends in different clinical scenarios but not why those recommendations emerge from specific feature values. This lack of quantitative feature importance metrics poses challenges for regulatory approval. Agencies like the FDA require explainable AI systems \citep{fda2021ai}, yet without methods like SHAP values, LEG saliency, or attention weights, clinicians cannot validate whether the policy's decision logic aligns with medical reasoning at a fine-grained level.

Our work directly addresses this gap by systematically comparing offline RL algorithms, including CQL which shares conceptual similarities with the AI Clinician's conservative fitted Q-iteration approach, on both survival outcomes and quantitative interpretability metrics. While Komorowski's study established that RL-derived policies can achieve high performance in sepsis treatment, our contribution demonstrates that algorithm choice critically affects interpretability. Conservative offline methods like CQL produce approximately 580-fold stronger feature importance signals than vanilla DQN, enabling clinicians to understand and validate treatment recommendations. This quantitative interpretability benchmark complements prior performance-focused work and provides actionable guidance for selecting algorithms suitable for clinical deployment. While CQL has been explored for sepsis treatment in recent work, these studies similarly did not evaluate policy interpretability using gradient-based saliency methods, motivating our systematic LEG-based comparison across offline and online RL paradigms within the gym-sepsis simulator.


% ============================================================
\subsection{Offline Reinforcement Learning}\label{sec:related:offline}

Offline RL learns from fixed datasets without environment interaction, addressing distributional shift when learned policies select out-of-distribution actions \citep{levine2020offline}. Behavior Cloning \citep{pomerleau1991bc} treats offline RL as supervised imitation learning, avoiding distributional shift but cannot improve beyond the behavioral policy. Conservative Q-Learning \citep{kumar2020cql} adds a conservatism penalty that discourages high Q-values for out-of-distribution actions, providing safety guarantees for healthcare. The CQL penalty encourages simpler Q-function representations aligned with the behavioral policy. When the behavioral policy follows interpretable threshold rules, CQL may learn Q-functions with strong gradients detectable by saliency analysis. Deep Q-Network \citep{mnih2015dqn} combines Q-learning with deep networks and experience replay. Originally designed for online learning, it can be adapted to offline settings but tends to overestimate out-of-distribution action values.


% ============================================================
\subsection{Interpretability Methods in RL}\label{sec:related:interp}

Regulatory agencies require explainable AI for medical decision support \citep{holzinger2017xai_healthcare}, yet deep RL policies are notoriously opaque. Greydanus et al. \citeyearpar{greydanus2018leg} introduced Linearly Estimated Gradients, a perturbation-based method that approximates Q-function gradients via local linear regression, producing saliency maps highlighting which features drive action selection. We adopt LEG because it is model-agnostic, enabling fair comparison across algorithms, produces quantitative saliency scores, and aligns with clinical intuition where clinicians weight physiological indicators. While prior sepsis RL work has explored interpretability through policy distillation and attention mechanisms, to our knowledge no study has systematically compared gradient-based saliency across multiple RL algorithms in the sepsis domain.


% ============================================================
\subsection{Research Gap}\label{sec:related:gap}

Despite a decade of sepsis RL research achieving strong retrospective performance, systematic quantitative interpretability evaluation using gradient-based saliency methods remains limited. While offline RL methods differ fundamentally in handling distributional shift, the relationship between these algorithmic differences and interpretability has not been thoroughly examined. Our work addresses this gap by jointly benchmarking offline and online RL on both survival outcomes and LEG-based interpretability within the gym-sepsis simulator, providing quantitative evidence on how algorithm choice affects the performance-interpretability trade-off.

% End of Related Work section